\item Un bloque metálico de $10.0 \text{ kg}$ que mide $12.0 \text{ cm}$ por $10.0 \text{ cm}$ por $10.0 \text{ cm}$, está suspendido de una balanza y sumergido en agua, como se muestra en la figura P14.11b. La dimensión de $12.0 \text{ cm}$ es vertical y la parte superior del bloque está $5.00 \text{ cm}$ debajo de la superficie del agua. (a) ¿Cuáles son las magnitudes de las fuerzas que actúan sobre las partes superior e inferior del bloque debidas al agua circundante? (b) ¿Cuál es la lectura en la balanza de resorte? (c) Demuestre que la fuerza de flotación es igual a la diferencia entre las fuerzas sobre las partes inferior y superior del bloque.
